\documentclass[11pt]{letter}
\usepackage[margin=1in]{geometry}
\usepackage{hyperref}

\signature{Bhaskar Gurram\\Zasti Inc., Ashburn, VA, USA\\gurrambhaskar.ai@gmail.com}
\address{Bhaskar Gurram\\Zasti Inc.\\Ashburn, VA 20147\\USA}

\begin{document}

\begin{letter}{Dr.\ Jim Jansen\\Editor-in-Chief\\Information Processing \& Management}

\opening{Dear Dr.\ Jansen,}

Thank you for the constructive feedback on our submission
``AgentHallu-Bench: Measuring and Mitigating Error Propagation in
Tool-Using Language Agents.''
We have made substantial revisions to address every point in your letter,
and we are pleased to resubmit the revised manuscript, now titled
\textbf{``Evaluating Tool-Using Language Agents: Judge Reliability,
Propagation Cascades, and Runtime Mitigation in AgentProp-Bench.''}

We renamed the benchmark from AgentHallu-Bench to AgentProp-Bench to
reflect its focus on propagation measurement and to avoid confusion with
the concurrently published ``AgentHallu'' paper (Liu et al., 2026,
arXiv:2601.06818), which pursues step-level attribution on a different
dataset.

Specific changes addressing your eight requests are detailed in the
accompanying Response to the Editor and summarized here:
(1)~the abstract now leads with specific numbers rather than background
(Abstract, ${\sim}$170 words);
(2)~explicit research objectives appear as a standalone Section~3;
(3)~a full dataset description is provided in Section~4 with domain
breakdown, difficulty distribution, and example tasks;
(4)~a dedicated Discussion and Implications section (Section~7) addresses
theoretical implications, practical guidance for agent designers, and
relevance to IR/IS evaluation methodology;
(5)~the manuscript has undergone thorough copy editing for accessibility
to non-native English speakers;
(6)~the paper has been expanded from approximately 5,200 to 9,500 words;
(7)~the literature review (Section~2) now includes 18 new 2024--2026
references from IR/IS and agent-hallucination venues, with at least
three from SIGIR, CIKM, IP\&M, and TOIS;
(8)~all references follow APA formatting via \texttt{elsarticle-num-names.bst}.

During revision, we also identified and corrected a methodological issue:
the original stage-indicator definitions coupled two stages, producing a
forced $r_{2,3}=1.00$. We corrected the indicators, re-judged all traces
with a validated 3-LLM ensemble, collected 100 stratified human labels
($\kappa=0.432$ vs.\ ensemble), and re-ran the runtime interceptor under
a concurrent no-interceptor control. All corrections are transparently
disclosed in the Response to the Editor (Section~9) and in Section~5.2
of the revised manuscript.

We believe the revised manuscript is a strong fit for IP\&M's audience,
as the primary contribution---demonstrating that substring-based evaluation
of agent outputs is at chance-level agreement with human annotation
($\kappa=0.049$), while a validated LLM-ensemble judge reaches moderate
agreement ($\kappa=0.432$)---has direct implications for evaluation
methodology in information retrieval and information science.

Per your request, we are willing to serve as reviewers for two manuscripts
submitted to IP\&M.

\closing{Best regards,}

\end{letter}
\end{document}
